\section{Interpreting Genetic Estimates and Policy Claims}
\label{sec:interpretation}
The preceding sections question identification and empirical support.
A separate issue remains even if a genetic decomposition is accepted:
what does it establish about social mechanisms, opportunity, and policy?
Keeping this conditional argument separate prevents an objection to the
estimates from substituting for an objection to their interpretation.

\subsection{Two Distinct Critiques}

\citet[pp.~344--346]{Goldberger1979} challenges the inference from a
variance decomposition to the effects and desirability of interventions.
The issue is what the parameter measures, even if it is known without
error. In the additive model with uncorrelated genetic and environmental
components, heritability is an $R^2$: the share of outcome variance
explained by the genetic component. It is not a response coefficient.
Goldberger explicitly distinguishes an explanatory variable's
contribution to $R^2$ from its regression slope. Policy analysis requires
causal responses to specified changes and their benefits and costs;
a variance share supplies neither. A small environmental variance component
can reflect little environmental variation rather than a small causal
response. Even if genetic effects operate entirely through biological
pathways, high heritability does not establish that those pathways cannot
be modified or their consequences compensated. Goldberger's eyeglasses
example illustrates why the genetic origin of a disadvantage does not
establish that remedying it is ineffective or prohibitively costly.

\citet[pp.~729--734]{Jencks1980} challenges the interpretation of genetic
and social causes as mutually exclusive. Genes can affect the environments
people select and the treatment they receive; those environments can then
affect outcomes. Thus a genetic variance component may include socially
mediated effects, and $1-h^2$ does not cap the causal importance of social
environments. His ``Genetic = Physical and Environmental = Social''
fallacy confuses the source of variation with the mechanism transmitting
its effects. Table~1 contrasts PKU (genetic/physical), lead poisoning
(nongenetic/physical), sexism (genetic differences eliciting social
treatment), and language spoken at home (nongenetic/social). Jencks also
qualifies these categories: lead exposure is socially patterned, while
PKU's consequences depend on diet, screening, and treatment. Entire
causal chains cannot be divided neatly into physical and social causes.

These are complementary emphases, not exclusive claims of intellectual
ownership: both authors discuss modifiability and policy. Neither critique
requires genetic estimates to be biased, and neither demonstrates that a
particular intervention works. The distinction is also developed in
\citet[secs.~3.4 and 3.6]{BenjaminEtAl2026}.


\subsection{Direct Genetic Effects Can Operate through Social Responses}
Clark describes the direct path as one in which ``the genome dictates
child social aptitudes, independent of any interaction with the parents''
\citep[p.~286]{Clark2026}. This excludes pathways that the modern direct-effect
estimand includes, as Section~\ref{sec:model-components} explains. An effect
of the child's own genotype can operate through parental attention,
teachers, or peers. Evidence against parental genetic nurture would not
establish that these responses are absent. Confounding asks whether an
effect has been identified; mediation asks how that effect operates.
Treating ``direct'' as narrowly physiological conflates the two.

\subsection{Clark's Qualifications and His Stronger Conclusions}
The book also contains explicit acknowledgments of these distinctions.
Clark writes: ``This does not in itself imply that social interventions
cannot change social outcomes'' \citep[p.~80]{Clark2026}, and recognizes
that intervention remains an empirical question even without convincing
evidence of genetic nurture \citep[p.~86]{Clark2026}. He notes that
heritability depends on the range of family environments, and allows
that replacement care could explain small effects of parental loss \citep[pp.~14,191,288]{Clark2026}.
The criticism cannot fairly be that he never recognizes such possibilities.

The problem is their inconsistent application. The claim that direct
genetic transmission implies excessive education spending (p.~9), or
that social outcomes can change only when genetic endowments change
(p.~10), does not follow from a variance decomposition \citep[pp.~9--10]{Clark2026}. Nor does a
genetic explanation establish meritocracy \citep[p.~7]{Clark2026}: discriminatory social
responses to inherited characteristics can themselves generate genetic
associations or effects. Appendix~\ref{app:quotations} preserves ten
critical passages alongside ten qualifications, with their contexts.
These are tensions in the interpretation, not proof that every empirical
claim in the book is false.

\subsection{Education Spending: Effects, Costs, and the Policy Objective}
Clark calls excessive education expenditure an implication of direct
genetic transmission and points to persistent mobility despite narrowing
gaps in schooling \citep[p.~9]{Clark2026}. But whether expenditure is
excessive requires marginal benefits and costs. The relevant benefits can
include absolute earnings, health, and consumption as well as changes in
inequality or relative mobility. A program that raises everyone's reading
ability can be valuable even if the same families retain their relative
positions. Unchanged ranks or intergenerational correlations therefore do
not establish that an intervention was ineffective or not worth its cost.

The intervention-specific evidence discussed in
Section~\ref{sec:environmental-evidence} is relevant to this question, but
its estimands must be matched to the policy claim. A local effect of a
school-leaving reform, an average return to schooling, and the return to
an expansion of public expenditure need not coincide. The genetic model
cannot supply the missing welfare comparison.

\subsection{Historical and Population Extrapolations Need Additional Evidence}
The same logic applies to claims about future outcomes and population
composition. Genetic contributions estimated under current institutions
do not establish invariant effects under different schools, labor markets,
or tax systems. Clark's acknowledgment that social and legal restrictions
limited the prospects of Jewish migrants in their countries of origin
\citep[p.~291, n.~105]{Clark2026} illustrates the relevance of those
conditions. Conclusions about future fiscal costs or historical innovation
require their own evidence; they cannot be read directly from a fitted
kinship model. Neither the identification critique nor the interpretive
critique implies that any particular policy must succeed. They establish
why success or failure must be assessed on evidence appropriate to it.
